\documentclass[tikz,border=2mm]{standalone}
\usepackage{tikz}
\usetikzlibrary{positioning, arrows.meta, shapes.geometric, fit, calc}
\usepackage{amsmath,amssymb,bm}

\begin{document}

% More professional, academic color scheme for mathematical physics journal
\definecolor{enccolor}{HTML}{34495E}      % Dark blue-gray
\definecolor{scalecolor}{HTML}{7F8C8D}    % Medium gray
\definecolor{latcolor}{HTML}{E67E22}      % Muted orange
\definecolor{deccolor}{HTML}{2C3E50}      % Dark slate
\definecolor{auxcolor}{HTML}{95A5A6}      % Light gray
\definecolor{graytext}{HTML}{2C3E50}      % Dark text
\definecolor{lightgray}{HTML}{BDC3C7}     % Light gray

\begin{tikzpicture}[
    node distance=1.2cm and 1.8cm,
    >={Stealth[length=2.5mm]},
    block/.style={
        rectangle,
        rounded corners=3pt,
        minimum width=2.0cm,
        minimum height=1.1cm,
        text centered,
        draw=white,
        line width=1.2pt,
        font=\sffamily\small
    },
    smallblock/.style={
        rectangle,
        rounded corners=2pt,
        minimum width=1.3cm,
        minimum height=0.7cm,
        text centered,
        draw=white,
        line width=1pt,
        font=\sffamily\footnotesize
    },
    circlenode/.style={
        circle,
        minimum size=0.9cm,
        text centered,
        draw=white,
        line width=1.2pt,
        font=\sffamily\bfseries
    },
    arr/.style={
        ->,
        line width=1pt,
        color=graytext
    },
    label/.style={
        font=\sffamily\bfseries\footnotesize
    }
]

% === INPUTS ===
\node[font=\bfseries, text=graytext] (input_x) at (0, 2.5) {$\mathbf{X}$};
\node[below=0.08cm of input_x, font=\tiny, text=lightgray] {$(N, 7, 65)$};

\node[font=\bfseries, text=graytext] (input_m) at (0, -2.5) {$\mathbf{m}$};
\node[below=0.08cm of input_m, font=\tiny, text=lightgray] {$(N, 7)$};

% === SHAPE ENCODER ===
\node[block, fill=enccolor, text=white, right=2.0cm of input_x] (shape_enc) {Shape Enc.\\[-2pt]{\tiny Bi-LSTM}};

% === SCALE ENCODER ===
\node[block, fill=scalecolor, text=white, right=2.0cm of input_m] (scale_enc) {Scale Enc.\\[-2pt]{\tiny MLP}};

% === Z_SHAPE / Z_SCALE ===
\node[smallblock, fill=enccolor!80, text=white, right=1.8cm of shape_enc] (z_shape) {$\mathbf{z}_{\text{sh}}$};
\node[smallblock, fill=scalecolor!80, text=white, right=1.8cm of scale_enc] (z_scale) {$\mathbf{z}_{\text{sc}}$};

% === CONCATENATION NODE ===
\node[circlenode, fill=lightgray, text=white, right=2.2cm of $(z_shape)!0.5!(z_scale)$] (concat) {$\oplus$};

% === MU/SIGMA ===
\node[smallblock, fill=latcolor, text=white, above right=0.4cm and 1.5cm of concat] (mu) {$\bm{\mu}$};
\node[smallblock, fill=latcolor, text=white, below right=0.4cm and 1.5cm of concat] (sigma) {$\bm{\sigma}$};

% === LATENT Z ===
\node[circlenode, fill=latcolor, text=white, right=3.0cm of concat] (z) {$\mathbf{z}$};

% === DECODER ===
\node[block, fill=deccolor, text=white, right=2.2cm of z] (decoder) {Decoder\\[-2pt]{\tiny LSTM}};

% === OUTPUT X_HAT ===
\node[font=\bfseries, text=graytext, right=1.8cm of decoder] (output_x) {$\hat{\mathbf{X}}$};
\node[below=0.1cm of output_x, font=\tiny, text=lightgray] {$(N, 7, 65)$};

% === SCALE PREDICTOR ===
\node[block, fill=auxcolor, text=white, below=2.2cm of decoder, minimum height=0.9cm] (mlp) {Predictor\\[-2pt]{\tiny MLP}};
\node[font=\bfseries, text=graytext, right=1.8cm of mlp] (output_m) {$\hat{\mathbf{m}}$};
\node[below=0.1cm of output_m, font=\tiny, text=lightgray] {$(N, 7)$};

% === ARROWS ===
% Inputs to encoders
\draw[arr] (input_x) -- (shape_enc);
\draw[arr] (input_m) -- (scale_enc);

% Encoders to latent representations
\draw[arr] (shape_enc) -- (z_shape);
\draw[arr] (scale_enc) -- (z_scale);

% Latent representations to concatenation
\draw[arr] (z_shape.east) -- ++(0.3,0) |- (concat);
\draw[arr] (z_scale.east) -- ++(0.3,0) |- (concat);

% Concatenation to mu/sigma
\draw[arr] (concat.east) -- ++(0.3,0) |- (mu.west);
\draw[arr] (concat.east) -- ++(0.3,0) |- (sigma.west);

% Mu/sigma to z (sampling) - with better label
\draw[arr] (mu.east) -- ++(0.3,0) |- (z);
\draw[arr] (sigma.east) -- ++(0.3,0) |- (z);

% Z to decoder
\draw[arr] (z) -- (decoder);

% Decoder to output
\draw[arr] (decoder) -- (output_x);

% Z to predictor
\draw[arr] (z.south) -- ++(0,-0.5) -| (mlp.north);

% Predictor to scale output
\draw[arr] (mlp) -- (output_m);

% === ANNOTATIONS ===
% Reparameterization sampling label (better placement)
\node[below=0.25cm of z, font=\scriptsize, text=latcolor!70, align=center] {$\mathcal{N}(\bm{\mu}, \bm{\sigma}^2)$};

% Training only label
\node[above=0.15cm of input_m, font=\tiny, text=scalecolor!80, align=center] {\textit{(training)}};

% === SECTION LABELS (Better positioning) ===
\node[label, text=enccolor, above=0.6cm of shape_enc] {Shape Encoding};
\node[label, text=scalecolor, above=0.6cm of scale_enc] {Scale Conditioning};

% Latent Space label centered over the latent region
\node[label, text=latcolor, above=0.7cm of $(mu)!0.5!(z)$] {Latent Space};

\node[label, text=deccolor, above=0.6cm of decoder] {Reconstruction};
\node[label, text=auxcolor, below=0.5cm of mlp] {Scale Prediction};

\end{tikzpicture}

\end{document}
